\chapter{Decision under Uncertainty: Criteria, Confidence and Sources}
\label{sec:p17-rationality-8-4}

Hill's survey asks how uncertainty about evidence should enter rational
choice. It distinguishes a thesis about beliefs---one additive probability
per event---from a thesis about decisions---maximizing expected utility
under those beliefs. Its alternatives change either or both theses.
The survey's emphasis is normative, even where a model was originally
developed for descriptive accuracy or economic tractability
\cite[chapter~8.4, pp.~495--504]{Hill2021DecisionUncertainty}.
This account preserves its multiple-prior, nonadditive, confidence and
probability-based approaches, rather than treating them as names for maxmin.

\section{The source's decision domain and opening challenge}
\label{sec:p17-rat84-domain}

\begin{hbdefinition}{def:p17-rat84-acts}{States, consequences, acts and the utility scale}
Let \(S\) be a finite nonempty state space, \(X\) a consequence space,
and \(f:S\to X\) an act. An event \(E\) is a subset of \(S\).
Fix the source's common, state-independent utility \(u:X\to\mathbb R\)
and write \(v_f(s)=u(f(s))\). The symbols \(1,0,x\) denote
consequences of those utility values, not necessarily money amounts.
The bet \(1_E0\) gives utility one on \(E\) and zero otherwise;
more generally \(x_Ey\) uses the stated winning and losing utilities.

Write \(\Delta\) for the additive probability vectors on \(S\), and
retain the source's notation
\[
                  EU_p f=\sum_{s\in S}p(s)u(f(s)).
\]
A value \(V\) represents weak preference when
\(f\succeq g\) iff \(V(f)\geq V(g)\).
For calculations requiring arbitrary utility vectors, we explicitly use
the rich domain \(\mathbb R^S\), with all sure utility levels available.
Other comparisons need only the finite acts displayed.
\end{hbdefinition}

The source calls a decision \emph{under risk} when probabilities are
given, and \emph{under uncertainty} when they are not. Its unknown urn
contains black and white balls in an unknown proportion; its known urn
has half of each. This is not a claim of complete absence of information:
the possible colors and the known urn's proportion are specified.
Keynes's weight of evidence and Popper's ideal-evidence discussion motivate
the distinction between a best probability and the support for it.
Replacing direct knowledge of an urn's composition by sampling changes
the evidential route; a few sampled draws do not by themselves establish
an exactly known composition.
Hill notes that precise Bayesian beliefs can differ in their judgments
about composition or robustness to further observation; the belief-only
objection is therefore not decisive by itself.
\cite[Sec.~1 and notes~1--2, pp.~495--496, 503]{Hill2021DecisionUncertainty}

Let \(B_k,W_k,B_u,W_u\) denote the black/white events in the known
and unknown urns. In the source's state order the offered bets are
\[
\begin{array}{c|rrrr}
 &B_k\cap B_u&B_k\cap W_u&W_k\cap B_u&W_k\cap W_u\\\hline
 1_{B_k}0&1&1&0&0\\
 1_{B_u}0&1&0&1&0\\
 1_{W_k}0&0&0&1&1\\
 1_{W_u}0&0&1&0&1
\end{array}
\]
The Ellsberg pattern prefers the known-urn bet for either color. The
contrast is already mathematically impossible under one common additive
prior and the same favorable prize pair: \(p(B_k)>p(B_u)\) would
require \(p(W_k)<p(W_u)\). The complete original comparison is
\cref{thm:p17-ellsberg-two-urns}. Its hypothetical contradiction is not
a population-frequency estimate. Hill's cited reports of ambiguity
aversion motivate alternatives; they do not establish that such preferences
are uniquely rational. The combined challenge is that weight of evidence
has no additional role in ordinary one-prior expected-utility choice.

\section{Sets of priors: different comparisons and different choices}
\label{sec:p17-rat84-multiple-priors}

\begin{hbdefinition}{def:p17-rat84-prior-rules}{Maxmin, alpha-maxmin, unanimity and E-admissibility}
Supply a nonempty compact set \(\mathcal C\subseteq\Delta\). The source's
\emph{maximin expected utility}, called maxmin in the existing treatment,
and its fixed \(\alpha\in[0,1]\) generalization are
\[
\begin{aligned}
 V_M(f)&=\min_{p\in\mathcal C}EU_p f,\\
 V_\alpha(f)&=\alpha\min_{p\in\mathcal C}EU_p f
                   +(1-\alpha)\max_{p\in\mathcal C}EU_p f.
\end{aligned}
\]
The weighting convention puts \(\alpha=1\) at complete pessimism.
Unanimity instead defines a relation,
\[
       f\succeq_U g\quad\Longleftrightarrow\quad
       EU_p f\geq EU_p g\quad\text{for every }p\in\mathcal C.
\]
Its strict part requires the weak comparison and failure of the reverse;
using strict inequalities for every prior is a different policy.
For a specified nonempty finite menu \(\mathcal M\), an act is
\emph{E-admissible} if it maximizes \(EU_p\) on that menu for at
least one \(p\in\mathcal C\). This is menu-dependent and is not
the comparison of two independently minimized utilities.
\end{hbdefinition}

These are the chapter's equations~(1)--(3) and its E-admissibility contrast
\cite[Sec.~2, pp.~496--497]{Hill2021DecisionUncertainty}.
The source also mentions generalizations allowing \(\alpha\) to
depend on the act. They add a function \(\alpha(f)\) as an input;
the fixed-alpha properties here cannot be assumed for that larger class.
The common prior set does not imply that a single prior minimizes every
act. Nor does a rule selecting an act specify how the set was learned.
The corresponding finite method and its general lottery-act representation
are \cref{def:p17-finite-maxmin,thm:p17-mmeu-representation}.
Hill's fixed-utility deterministic-act survey is not that general
axiomatic domain.

\begin{hbtheorem}{thm:p17-rat84-criteria}{proved=annex:p17.rat84-criteria}{The source's urn values and a menu separating the rules}
For
\[
 \mathcal C^*=\{p\in\Delta:.3\leq p(B_u)\leq.7,
                                   \ p(B_k)=.5\},
\]
the maxmin values of the two known-urn bets are \(.5\), and those
of the two unknown-urn bets are \(.3\). Their alpha-maxmin values
are respectively \(.5\) and \(.7-.4\alpha\), so the source's
strict known-urn pattern holds iff \(\alpha>1/2\).
Unanimity cannot compare either known-urn bet with its same-color
unknown-urn counterpart under this set.

For the two-state menu of utility vectors
\(a=(1,0),b=(0,1),c=(.4,.4)\), with \(\mathcal C=\Delta\),
all three are undominated under the strict part of unanimity;
only \(a,b\) are E-admissible; maxmin uniquely chooses \(c\).
Alpha-maxmin chooses \(c\) for \(\alpha>.6\), chooses \(a,b\)
for \(\alpha<.6\), and ties all three at \(\alpha=.6\).
\end{hbtheorem}

The menu contrast makes the input distinction consequential. Unanimity
can withhold a ranking, E-admissibility asks whether an act is best under
one common prior for the entire menu, and maxmin completes the comparison
by a particular cautious scalar. Deferral, a status quo, random choice or
sharpening to one prior are additional mechanisms, not consequences of
an incomplete relation alone. Choosing a different favorable sharpening
for each transaction can also be incoherent; a completion or dynamic
policy must be specified consistently.

\subsection{Choice axioms and the scope of a representation theorem}
\label{sec:p17-rat84-choice-axioms}

Hill organizes normative assessment in two ways: examine the consequences
for choice, and examine the interpretation of the representing attitudes.
For the first, weak order (WO) means completeness and transitivity. The
sure-thing principle (STP) says that changing common consequences on
states where two acts agree leaves their comparison unchanged. In the
four-state table, changing the common first and last columns turns
the black-bet comparison into the reversed white-bet comparison, so
the Ellsberg pattern violates STP without estimating any prior.

Maxmin retains WO and weakens STP. Unanimity retains STP and transitivity
but drops completeness. A representation theorem states precise axioms
equivalent to existence of a representing model, with exact uniqueness
qualifications. The chapter's figure~8.4.1 is a schema for such a theorem;
the actual Gilboa--Schmeidler axioms, proof and nondegeneracy qualification
are supplied in \cref{def:p17-mmeu-domain,thm:p17-mmeu-representation}.
A few urn choices neither identify a unique prior set nor prove that
the axioms govern the agent's full preferences.

Two further source properties are P4, stake independence for comparisons
of event bets using the same winning and losing consequences, and UA,
the following uncertainty-aversion hedge:
\[
 E\cap F=\varnothing,\quad 1_E0\sim1_F0
 \quad\Longrightarrow\quad (1/2)_{E\cup F}0\succeq1_E0.
\]
The half refers to utility, not half the money under an arbitrary
nonlinear monetary utility. Maxmin's concavity establishes this hedge;
fixed alpha-maxmin can drop it. Its P4 property remains, however.
Hill also cites a version with ten possible colors in which choosing
the unknown urn can be attractive. This is the survey's brief motivating
example of a departure from universal ambiguity aversion, not a reported
population frequency. Its cited Wakker treatment is not substituted by
the natural-event studies below; no uninspected ten-color sample count
or effect size is supplied.
\cite[Sec.~2.1 and notes~5--7, pp.~497--498, 503]{Hill2021DecisionUncertainty}

\begin{hblemma}{thm:p17-rat84-axioms}{proved=annex:p17.rat84-axioms}{P4, the hedge and the static/dynamic boundary}
Maxmin and fixed alpha-maxmin satisfy P4 for positive utility prize
differences. Unanimity satisfies STP. Maxmin satisfies UA; with
two point priors and equal complementary bets, alpha-maxmin violates
UA when \(\alpha<1/2\).

Static completeness and transitivity do not alone establish dynamic
consistency under conditioning. On four states retain
\(p=(1/2,0,1/2,0)\), \(q=(0,1/2,0,1/2)\).
Maxmin prefers \(f=(1,0,0,1)\) to the sure act \(g=.4\) before
learning whether \(E=\{1,2\}\) occurs. After conditioning both
priors on either \(E\) or \(E^c\), it prefers \(g\) to \(f\).
\end{hblemma}

The source describes dynamic objections as a central dispute, citing
arguments for and against retaining WO or STP. Their assessment requires
the actual information tree, update rule and feasible actions. A static
representation is not a blanket guarantee of immunity to every sequence
of conditional decisions or sure-loss argument. The counterexample
identifies a conditional-plan inconsistency, not by itself a money pump
or a verdict on every ambiguity model. P15's natural and regular
conditioning interfaces likewise do not select a complete decision policy.

\subsection{Separating evidence from willingness to bear uncertainty}
\label{sec:p17-rat84-belief-taste}

With \(\mathcal C_0=\{p:0\leq p(B_u)\leq1,p(B_k)=.5\}\),
maxmin values the unknown-urn black bet at zero, below a known bet
with a one-in-a-hundred winning chance. Using the source's narrower
unknown-urn range \([.4,.6]\), while retaining the known-urn information
for joint comparisons, gives value \(.4\) instead of zero.
This alone does not explain why the range is narrower: it may express
additional beliefs, an insufficient-reason judgment, or less willingness
to act on extreme possibilities. Objective information and the set
revealed by preferences need not coincide.

\begin{hbdefinition}{def:p17-rat84-comparative-aversion}{Comparative uncertainty aversion}
For agents with the same utility scale, agent 1 is at least as
uncertainty averse as agent 2 when, for every act \(f\) and sure
utility \(z\), \(f\succeq_1 z\) implies \(f\succeq_2 z\).
An uncertain act acceptable to the more cautious agent must also
be acceptable to the less cautious one. The common scale and the
available sure consequences are part of this definition.
\end{hbdefinition}

\begin{hbtheorem}{thm:p17-rat84-comparative}{proved=annex:p17.rat84-comparative}{The finite canonical-set version of figure 8.4.2}
On \(\mathbb R^S\) with all sure utility levels available, let
agents 1 and 2 use maxmin with nonempty compact convex prior sets
\(\mathcal C_1,\mathcal C_2\). Agent 1 is at least as uncertainty
averse as agent 2 exactly when
\(\mathcal C_2\subseteq\mathcal C_1\).
For nonconvex representing lists the statement concerns their closed
convex hulls. Restricted observations or a sparse consequence domain
do not by themselves imply this converse.
\end{hbtheorem}

This is Hill's interpretational point: enlarging a representing prior
set corresponds to greater caution, which is commonly treated as a taste,
yet the same set is often called a belief representation. The theorem
does not show that beliefs and values must be mixed; it shows why
choice-based recovery alone cannot decide which interpretation is intended.
The distinction matters when experts report factual uncertainty and
others choose values or when learning is to be separated from decision.
Hill extends the concern to unanimity and capacities, and notes that
alpha-maxmin's representing \(\alpha\) and prior set can trade off.
Neither a revealed parameter nor a label ``credal set'' establishes
that it encodes evidence only.
\cite[Secs.~2.2--2.3 and note~9, pp.~498--500, 503]{Hill2021DecisionUncertainty}

\section{Capacities and Choquet evaluation}
\label{sec:p17-rat84-capacity}

\begin{hbdefinition}{def:p17-rat84-choquet}{A capacity and its finite Choquet value}
A normalized capacity \(\nu:2^S\to[0,1]\) has
\(\nu(\varnothing)=0,\nu(S)=1\), and is monotone under inclusion;
additivity is not required. Sort the distinct values of a utility
vector as \(x_1>\cdots>x_m\), and put
\(A_i=\{s:v_f(s)\geq x_i\}\). The finite Choquet value is
\[
        V_\nu(f)=x_m+\sum_{i=1}^{m-1}(x_i-x_{i+1})\nu(A_i).
\]
The lowest-value baseline makes the formula valid for signed utilities.
For nonnegative values it is the source's equation~(4), with the
last following value set to zero. For a constant vector it returns
that constant.
A capacity is \emph{convex}, in the source's terminology, when
\(\nu(E\cup F)+\nu(E\cap F)\geq\nu(E)+\nu(F)\).
Its core is
\(\{p\in\Delta:p(E)\geq\nu(E)\text{ for every }E\subseteq S\}\).
\end{hbdefinition}

\begin{hbtheorem}{thm:p17-rat84-choquet}{proved=annex:p17.rat84-choquet}{Dominance, event weights and the convex-core identity}
The Choquet value is monotone in the state utility vector, is invariant
under the sorting of tied values, and satisfies
\(V_\nu(x_Ey)=y+(x-y)\nu(E)\) for \(x>y\), hence P4.
For a convex capacity its core is nonempty, and
\[
              V_\nu(f)=\min_{p\in\operatorname{core}(\nu)}EU_p f.
\]
Consequently this case satisfies UA. A general capacity need not:
on two states with singleton capacities \(.75,.75\), each
complementary bet has value \(.75\), while the half hedge has
value \(.5\).

Naively substituting capacities for probabilities of each distinct
\emph{outcome event} can violate dominance. With singleton capacities
\(.1,.1\), that substitution values \((1,1)\) at 1 but the
strictly better \((1.1,1.2)\) at \(.23\). Choquet instead values
the better act at \(1.11\). Substituting singleton capacities on
a fixed state partition is another operation: it is monotone with
nonnegative weights, but need not normalize constants or represent
the capacity on nonsingleton events.
\end{hbtheorem}

This preserves the source's reason for changing the aggregation formula,
without attributing a dominance reversal to a fixed nonnegative weighted
sum that cannot have one. Convex capacities, including belief-function
examples, can yield the same criterion as maxmin over their core. That
does not make every arbitrary prior-set lower envelope a convex capacity
or make all capacities additive. The event bounds and full gamble
information are distinguished in \cref{sec:p15-imprecise-probability}.
The general Schmeidler axiomatization remains a cited boundary of the
survey; the finite core identity used here is fully justified locally.
\cite[Sec.~3, p.~500]{Hill2021DecisionUncertainty}

\section{Confidence levels, penalties and a second-order structure}
\label{sec:p17-rat84-confidence}

Hill's third family adds structure over \(\Delta\): an ordering of
probability judgments, or equivalently a nested family \(\Xi\) of
prior sets. The judgments valid under all priors in a larger set are
fewer and therefore can be held with higher confidence. His figure~8.4.3
has, from low to high confidence, a single \(p_0\) with both urn
marginals \(.5\); a set with unknown marginal \([.3,.7]\); and
a set with unknown marginal \([.1,.9]\). All retain known marginal
\(.5\). This ordering is not a probability distribution over priors.

An additional function \(D\) chooses the confidence set appropriate
to the decision. Its maxmin version is
\[
                       V_D(f)=\min_{p\in D(f)}EU_p f.
\]
The source motivates higher confidence requirements for larger stakes.
Thus the family's belief structure \(\Xi\) and the choice of how
much confidence to require can be separately modeled. Hill cites
representation results for the restricted confidence models, not for
every arbitrary act-dependent selector.
\cite[Sec.~4, pp.~500--502]{Hill2021DecisionUncertainty}

\begin{hblemma}{thm:p17-rat84-confidence}{proved=annex:p17.rat84-confidence}{A stakes contrast and why the selector needs restrictions}
With the source's nested sets, using \(p_0\) for unit-stake bets
ties known and unknown black bets at \(.5\). Requiring the high
confidence set for bets paying utility 100 values them at 50 and 10.
Thus a confidence-sensitive procedure can violate P4 while holding
its underlying nested family fixed.
An unrestricted selector can even violate statewise dominance:
on two states, evaluate \(f=(2,1)\) over all priors, and
\(g=(1.9,.9)\) under \((.5,.5)\). Then \(f\) is better in
every state but has value 1, below \(g\)'s 1.4.
\end{hblemma}

The latter is a boundary case for the schematic expression, not a
claim that Hill's fully axiomatized confidence theories allow arbitrary
selectors. Their admissible relationship between acts, stakes and
confidence levels matters. The comparison also differs from the existing
fixed-weight \(\rho\) security/best-estimate blend in P17: selecting
one member of a nested family is not mixing one common prior with
one fixed security set.

The chapter's two cardinal second-order criteria are
\[
       \min_{p\in\Delta}(EU_p f+c(p)),\qquad
       \min_{p\in\Delta}\frac{EU_p f}{\varphi(p)}.
\]
Here \(c\) is a penalty function and \(\varphi:\Delta\to[0,1]\)
is a fuzzy confidence function. Neither is a probability mass or a
decision weight on events. The first appears in variational preferences
and robust macroeconomic modeling; the second is the confidence-preference
approach. Both use cardinal structure, unlike a purely ordinal family.
The reciprocal needs a convention or exclusion at \(\varphi(p)=0\).
The bounded calculations below use positive confidence weights and
nonnegative utility; negative constants can otherwise interact with
the reciprocal in unintended ways.

\begin{hblemma}{thm:p17-rat84-cardinal}{proved=annex:p17.rat84-cardinal}{Two prior penalties and reciprocal weights are not one input}
For candidate priors \(p_1,p_2\) with expected utilities \((0,1)\),
penalties \((.2,0)\) give variational value \(.2\); for the
sure utility \(.3\), the value is \(.3\). Reciprocal confidence
weights \((.5,1)\) give values 0 and \(.3\), respectively.
For expected utilities \((.4,.6)\), the same penalty criterion
gives \(.6\), but the reciprocal criterion gives \(.6\),
above the raw maxmin \(.4\); the first agreement does not identify
their structures. A sharper contrast, \((.1,.6)\), gives
penalty value \(.3\) and reciprocal value \(.2\).

With a nonnegative lower-semicontinuous penalty on compact \(\Delta\),
having minimum zero, the minimum exists and constant acts retain their
utility. Maxmin over a closed \(\mathcal C\) is recovered by the
extended penalty zero on \(\mathcal C\), infinity outside; the
extended value is an explicit additional convention, not an ordinary
real-valued penalty. A positive continuous \(\varphi\), with
maximum one, likewise attains the reciprocal minimum and preserves
nonnegative constants. Its negative constants need not be preserved.
\end{hblemma}

The chapter cautions that a choice-based penalty or confidence weight
can also reflect uncertainty attitude. Labeling such functions as
confidence does not by itself separate evidence from tastes.
The penalty's units must track the fixed utility scale: under
\(u\mapsto au+b\), \(a>0\), also changing \(c\mapsto ac\)
preserves its ranking. The reciprocal model especially requires its
chosen utility zero to be respected.
\cite[Sec.~4, pp.~501--502]{Hill2021DecisionUncertainty}

\section{Probabilities retained: weighting, sources and smooth ambiguity}
\label{sec:p17-rat84-smooth}

A different family retains a precise \(p\) but deforms event weights
to \(w(p(E))\). If \(w\) is monotone with \(w(0)=0,w(1)=1\),
this gives a capacity and is evaluated by Choquet, not by blindly
replacing atomic masses. The same \(w\) cannot distinguish equal
probability same-prize bets across sources. Source dependence introduces
source-specific weights; a full model must also specify their application
to mixed-source acts. It does not follow from an event's probability alone.

A second-order probability \(\mu\) over priors is another input. On
finite support \(p_1,\ldots,p_m\), its linear reduction and the
smooth-ambiguity criterion are
\[
 \sum_j\mu_j EU_{p_j}f,
 \qquad V_S(f)=\sum_j\mu_j\varphi(EU_{p_j}f).
\]
In this last expression, \(\varphi\) has domain \(\mathbb R\)
and transforms expected \emph{utility}. It is not the preceding
fuzzy confidence function on \(\Delta\); the repeated glyph
has these different, explicitly typed roles in the source.
The source's integrals are these sums in the finite-support case.

\begin{hbtheorem}{thm:p17-rat84-reduction}{proved=annex:p17.rat84-reduction}{Linear reduction, source dependence and the smooth contrast}
Linear second-order expected utility equals one-prior expected utility
under \(\bar p=\sum_j\mu_jp_j\). With one weighting function,
two equal-probability event bets with the same prizes have equal
Choquet value. Neither construction alone yields the Ellsberg strict
two-color pattern at equal \(.5\) probabilities.

Let \(\mu\) put half its mass on a prior making the unknown-urn
black probability zero and half on one making it one; both assign
known-urn black probability \(.5\). With \(\varphi(x)=\sqrt{x}\)
on nonnegative utility, the known bet has smooth value
\(\sqrt{.5}\) and the unknown bet value \(.5\), for either color.
With \(\varphi(x)=x^2\), the ordering reverses; with affine
\(\varphi\), it disappears. More generally a strictly concave
increasing transform prefers the common expected utility to a
nonconstant distribution of expected utilities with that same mean.
\end{hbtheorem}

This separates first-order payoff uncertainty from uncertainty about
which probability model to use. \(\mu\) and the transform are distinct
ingredients; nonlinearity alone does not dictate the direction or
guarantee ambiguity aversion. Hill discusses the claimed clean separation
of beliefs and attitudes in this approach, but also its normative objection:
using an attitude parameter to make evidence weight matter may leave
unclear whether that weight is a characteristic of beliefs or tastes.
That objection is not settled by fitting a transform to choices.
\cite[Sec.~5 and note~10, pp.~502--503]{Hill2021DecisionUncertainty}

\section{Bounded empirical evidence and its conflict}
\label{sec:p17-rat84-evidence}

The source's normative survey refers to experimental tendencies, not a
single representative population estimate. The following primary set
is fixed to its source-dependence issue; observations are not promoted
to axioms about what everyone ought to choose.

\paragraph{Tversky and Fox's source comparisons.}
Their three studies used repeated choices between prospects and sure
cash amounts; certainty equivalents were inferred from switch intervals,
not freely reported. Study~1 analyzed 27 male Stanford basketball fans,
median age 21, paid \$15 for two sessions. Study~2 analyzed 40 male
football fans, median 21, with the same session payment and a chance of
playing a selected choice for real money. Study~3 analyzed 45 Stanford
psychology students, 28 male and 17 female, median 20, for course credit.
Additional participants excluded for internal inconsistency numbered
one, four and three respectively; a small unquantified number of outlying
responses was also excluded. Source blocks included known chip-draw
chances, sports, temperatures and a stock-index change. Chance blocks
preceded the natural-source blocks; the two temperature-source orders
in Study~3 were counterbalanced.
\cite[pp.~272--274]{TverskyFox1995Weighing}

Their source-preference index \(k\) counted concordant preferences
for an event and its complement across two sources, normalized by
patterns supporting either source. The reported median index for
basketball over chance was \(.76\), \(p<.05\); for Super Bowl
over chance it was \(.59\) but not statistically significant;
for San Francisco temperature over chance it was \(.76\), \(p<.01\).
The San Francisco/Beijing comparison was \(.86\), \(p<.01\).
These concern recruited competence-related groups and their supplied
tasks, not universal ambiguity seeking. Fitted decision weights also
used a utility exponent imported from an earlier study, with sensitivity
checks, rather than measuring it independently from every simple
prospect. The selected report supplies significance tests but not a
numerical confidence interval for these median indexes.
\cite[pp.~274, 280--282]{TverskyFox1995Weighing}

\paragraph{The eight-color source-function experiment.}
Abdellaoui and colleagues report 67 students comparing an eight-ball
urn with one of each color against an eight-ball urn of unknown
composition. Each faced 26 series of choices against ascending sure
payments; the switch midpoint was used as the certainty equivalent.
A randomly selected choice was played for real. Outcomes ranged from
EUR~0 to EUR~25. A power utility and source weighting were fitted from
the choices, and source functions at probabilities \(j/8\) were
then computed conditional on the utility estimate. One participant was
removed after the individual uniformity test, leaving 66 in the main
source-function analysis. The authors' August~2011 appendix identifies
recruitment from two French engineering schools by posters and online
registration, with 20 female and 47 male participants acquainted with
probability but not decision theory. Individual computer-assisted
interviews lasted about 20 minutes, and known/unknown urn order was
counterbalanced. The appendix lists 32 prospects; this count is not
silently identified with the main paper's 26 choice series.
Age distribution and recruitment dates remain unspecified here.
\cite[Sec.~III, pp.~703--707]{AbdellaouiEtAl2011Sources}
\cite[Secs.~A.2--A.3, pp.~4--7]{AbdellaouiEtAl2011SourcesAppendix}

For the seven-color event, median source weights were \(.94\) for
the known urn and \(.82\) for the unknown, with interquartile ranges
\([.80,.99]\) and \([.63,.93]\); the table reports a difference
test rounded to \(p=.000\), meaning rounded reporting, not literally
zero probability. At one color both medians were \(.19\), and the
between-source difference was not significant, \(p=.339\).
They report less sensitivity to likelihood and more pessimism for
the unknown source, with substantial individual heterogeneity. The
utility equality test \(p=.902\) does not establish exact equality
of the true utilities. Interquartile ranges describe the distribution
of estimated individual weights, not confidence intervals for a
population effect.

\paragraph{The opposing primary reanalysis.}
Harrison's May~2011 working-paper comment analyzes the same experiment,
not a new sample. It estimates power utility and Prelec weighting jointly
by maximum likelihood with a probabilistic choice-error specification,
including contextual normalization. He argues that conditioning later
weight estimates on first-stage utility point estimates ignores their
uncertainty. Of 66 participants, 13 could not be estimated stably and
were dropped. He reports that only four of the retained individuals
rejected source independence jointly for the weighting parameters at
the displayed \(.10\) threshold. His median utility-exponent estimate
was \(1.003\), with median standard error \(.67\), illustrating
substantial parameter uncertainty.
\cite[pp.~3--9, table~1]{Harrison2011SourcesComment}

The comment's informal descriptions of color/task counts do not replace
the original paper's protocol. The substantive conflict concerns inference
from estimated parameters and the error model. The original authors
also report alternative probabilistic-choice/maximum-likelihood analyses
in their web appendix as agreeing with their conclusions. That appendix
is itself part of the dispute: Harrison says the original maximum-
likelihood check was pooled, not the individual-level check he applies.
\cite[pp.~704--705]{AbdellaouiEtAl2011Sources}
\cite[p.~9]{Harrison2011SourcesComment}
The retained August~2011 appendix indeed calls its likelihood analysis
a representative-agent analysis, with a Fechnerian error term and
within-individual correlation correction. It reports an insensitivity
parameter of \(.81\) for known and \(.64\) for unknown probability,
with a significant difference, but no numerical standard errors in
that paragraph. The appendix postdates the May comment. Its aggregate
check and Harrison's individual parameter tests concern different
levels; failure to reject individual equality does not prove exact
source independence, nor does aggregate significance establish each
person's source effect. No participant-level recomputation is reported.
\cite[Sec.~A.5, p.~12]{AbdellaouiEtAl2011SourcesAppendix}
Thus this treatment records competing analyses and their
limitations, not a resolved verdict that source dependence is established
or refuted. Numerical instability exclusions, maintained utility forms,
within-person repeated observations and recruitment limits remain
consequential. None of the findings chooses a uniquely rational ambiguity
criterion.

\section{Relations and consultation cases}
\label{sec:p17-rat84-consultation}

The current P17 maxmin treatment fixes utility and priors or derives a
canonical model from precise lottery-act axioms. Hill asks whether such
a representing set is evidence, attitude, or both, and places it alongside
rules using capacities, nested sets, penalties or second-order probability.
The P15 lower-prevision interface can compute \(\underline E(v_f)\)
or \(\underline E(v_f-v_g)\); those answer different choice questions.
A coherent price model supplies neither a unique menu rule nor a
confidence threshold. Updating at lower probability zero, as distinguished
in \cref{thm:p15-ip-conditional-models}, is another separate issue.

\begin{hbmethod}{def:p17-rat84-use}{Select the criterion by its actual inputs and claim}
Fix states, consequences, the common utility and the feasible menu.
Record whether evidence is supplied as a precise prior, a set, event
capacities, a nested confidence family, a penalty, fuzzy prior confidence,
or a second-order probability. Add the separate criterion and attitude
parameters; do not infer them from the representation's name.
Compute all relevant expected utilities or ranked outcome sets, including
the source's limiting cases. For unanimous comparison retain incomparable
pairs; for E-admissibility retain the whole menu; for a scalar criterion
report all tied maximizers. If information will arrive, state the update
and continuation policy as well.

The three-act menu is a consultation failure case for substituting
maxmin for E-admissibility: maxmin chooses a compromise never best under
any retained prior. The convex-capacity theorem is a justified bridge
to maxmin, but its convexity hypothesis cannot be dropped. The nested
confidence example succeeds at making stakes affect evidential requirements,
while its unrestricted-selector counterexample warns against assuming
all such formulas respect dominance. The two-source empirical comparison
requires its actual elicitation, error model and population limits; a
reported descriptive tendency is not normative authorization.
\end{hbmethod}

Hill leaves the competition among rational criteria open. His notes locate
the connections to robust statistics, strict/weak unanimity, conditional
dilation, certainty independence and alternative second-order theories.
Those citations mark actual external boundaries; they do not allocate
the whole works of Savage, Walley, Wakker or the economic surveys here.
The argument is that uncertainty representation and choice policy must
be assessed together, while keeping evidence, outcome tastes and attitude
to uncertainty distinguishable wherever the model permits.
